\section{Theoretical Setup}\label{sec:setup}

\subsection{Problem Setting}

Fix an integer $N\geq 1$, a support radius $R>0$, a density cap
$\kappa\in(0,\infty)$, and a compact interval $\Theta\subset\mathbb R$.
Let
\[
b:\Theta\to\mathbb R,
\qquad
F=(F_1,\ldots,F_N):\Theta\to\mathbb R^N,
\qquad
\phi_\alpha(\theta)=b(\theta)+\langle\alpha,F(\theta)\rangle.
\]
Write $\|x\|_1=\sum_i|x_i|$, and write $|I|$ for the Lebesgue length of
an interval $I$.  Define
\[
A=(2R)^N\kappa.
\]
Let $\mathcal D_{N,R,\kappa}$ be the class of probability laws $\mu$ on
$\mathbb R^N$ having a Lebesgue density $f_\mu$ such that $f_\mu=0$
almost everywhere outside $[-R,R]^N$ and
$\|f_\mu\|_{L^\infty(\mathbb R^N)}\leq\kappa$.  No product structure or
coordinate independence is part of this definition.

The common one-dimensional Pfaffian representation used here has the form
\[
\eta_j'(\theta)
=P_j(\theta,\eta_1(\theta),\ldots,\eta_j(\theta)),
\qquad 1\leq j\leq q,
\]
with
\[
b(\theta)=Q_0(\theta,\eta(\theta)),
\qquad
F_i(\theta)=Q_i(\theta,\eta(\theta)),
\quad 1\leq i\leq N.
\]
In ambient parameter dimension $p=1$, we use the degree descriptors
\[
M=\max_{1\leq j\leq q}\deg P_j,
\qquad
\Delta_{\rm rnd}=\max_{1\leq i\leq N}\deg Q_i,
\qquad
\Delta_{\rm aff}=\max_{0\leq i\leq N}\deg Q_i.
\]
When $q=0$, the chain is absent and $M=0$.  The deterministic offset $b$
is counted in $\Delta_{\rm aff}$ but is not a random coefficient.

The root-feasible set is
\[
K_R=\{\theta\in\Theta:|b(\theta)|\leq R\|F(\theta)\|_1\}.
\]
For each coordinate $j$, let
$U_j=\{\theta\in\Theta:F_j(\theta)\neq0\}$ and define on $U_j$
\[
V_j(\theta)
=\left|\left(\frac b{F_j}\right)'(\theta)\right|
+R\sum_{i\neq j}
 \left|\left(\frac{F_i}{F_j}\right)'(\theta)\right|.
\]
For minimization, extend $V_j$ by $+\infty$ on $\Theta\setminus U_j$ and
set
\[
\Gamma_{\rm piv}(b,F;R)
=
\begin{cases}
\displaystyle
\sup_{\theta\in K_R}\min_{1\leq j\leq N}V_j(\theta),
&K_R\neq\varnothing,\\[1ex]
0,&K_R=\varnothing.
\end{cases}
\]
This is a static, pre-sampling conditioning functional.  Its finiteness is
a conclusion below, not an additional assumption.

For the adaptive sweep, let $j_*(\theta)$ be the least index attaining
$\min_jV_j(\theta)$ on $K_R$, and put
\[
E_j=\{\theta\in K_R:j_*(\theta)=j\},
\qquad
E_{j,m}=E_j\cap\{|F_j|\geq 1/m\},\quad m\geq1.
\]
For $\beta=\alpha_{-j}\in[-R,R]^{N-1}$, define on $E_j$
\[
T_j(\theta,\beta)
=-\frac{b(\theta)+\sum_{i\neq j}\beta_iF_i(\theta)}{F_j(\theta)}.
\]
The sweep applies the area formula to
$(\theta,\beta)\mapsto(\beta,T_j(\theta,\beta))$ in original coefficient
order, restricted by $|T_j|\leq R$, and then exhausts in $m$.  The sets
$E_j$ are disjoint in the parameter $\theta$, so the resulting estimate is
a sum of chart integrals rather than an $N$-fold union bound.

Two fixed families will be used.  For $\delta\in(0,1]$, the scale-stress
family on $[-1,1]$ is
\[
b_\delta(\theta)=0,
\qquad
F_\delta(\theta)=(1,\theta/\delta),
\qquad R=1.
\]
For an integer $d\geq1$, the exact monic family is
\[
b_d(\theta)=\theta^d,
\qquad
F_d(\theta)=(1,\theta,\ldots,\theta^{d-1}),
\qquad
p_\alpha(\theta)=\theta^d+\sum_{k=0}^{d-1}\alpha_k\theta^k.
\]
Only $\alpha=(\alpha_0,\ldots,\alpha_{d-1})\in[-R,R]^d$ is random.
Write $V_{\rm const}$ for the constant-coordinate pivot and
$V_{\rm top}$ for the $\theta^{d-1}$ pivot; these are the same coordinate
when $d=1$.  This family has
\[
q=0,
\quad M=0,
\quad N=d,
\quad \Delta_{\rm rnd}=d-1,
\quad \Delta_{\rm aff}=d,
\quad A=(2R)^d\kappa.
\]

\subsection{Assumptions}

\begin{assumption}[Shared one-dimensional Pfaffian regularity]
\label{assump:shared-pfaffian-chain}
The functions $b,F_1,\ldots,F_N$ are $C^1$ on $\Theta$ and have the common
triangular Pfaffian representation and the
$(q,M,\Delta_{\rm rnd},\Delta_{\rm aff},p=1)$ degree convention specified
above.
\end{assumption}

\begin{assumption}[Primitive no-forced-root nondegeneracy]
\label{assump:no-forced-root}
For every $\theta\in\Theta$,
\[
(b(\theta),F(\theta))\neq(0,0).
\]
Equivalently, no parameter value is a root for every coefficient vector.
Points with $F(\theta)=0$ and $b(\theta)\neq0$ remain admissible and are
root-free.  In particular, $F(\theta)\neq0$ for every $\theta\in K_R$.
\end{assumption}

\begin{assumption}[Arbitrarily correlated bounded joint density]
\label{assump:joint-density-cap}
The coefficient law is any $\mu\in\mathcal D_{N,R,\kappa}$.  Thus the cap
is on the full joint density and imposes no coordinate independence.
\end{assumption}

\subsection{Formalized Goal}

The objective is to prove, under
Assumptions~\ref{assump:shared-pfaffian-chain}--\ref{assump:joint-density-cap},
that $\Gamma_{\rm piv}(b,F;R)<\infty$ and that every
$\mu\in\mathcal D_{N,R,\kappa}$ and every interval $I\subseteq\Theta$ with
$|I|>0$ satisfy
\[
\Pr_{\alpha\sim\mu}\!\left[
  \exists\theta\in I:\phi_\alpha(\theta)=0
\right]
\leq
\kappa(2R)^{N-1}\Gamma_{\rm piv}(b,F;R)|I|
=\frac{A\Gamma_{\rm piv}(b,F;R)}{2R}|I|.
\]
The same constant must bound the law- and interval-uniform
probability-to-length ratio.  The scale-stress family must have
$\Gamma_{\rm piv}(b_\delta,F_\delta;1)=1/\delta$.  The monic family must
satisfy
\[
\Gamma_{\rm piv}(b_d,F_d;R)
\leq d+\frac{Rd(d-1)}2
\]
and, for every bounded interval $I\subset\mathbb R$ and every admissible law
on the lower coefficients,
\[
\Pr_{\alpha\sim\mu}\!\left[
  \exists\theta\in I:p_\alpha(\theta)=0
\right]
\leq
\kappa(2R)^{d-1}
\left(d+\frac{Rd(d-1)}2\right)|I|.
\]
These are ordinary-probability bounds, uniform over the stated laws and
intervals, with no hidden constant, horizon, or asymptotic qualification.
The exposed quantities are $N,R,\kappa,A,\Gamma_{\rm piv}$, and $|I|$;
$q,M,\Delta_{\rm rnd},\Delta_{\rm aff}$ are fixed family descriptors and
do not enter as an additional factor.  The norm modes are the $\ell_1$
feasibility test in $K_R$ and the coordinate-ratio variation in
$\Gamma_{\rm piv}$.

This is a material-partial objective.  It proves a fixed-family theorem and
the exact monic baseline, while leaving unresolved any polynomial
general-instance control of $\Gamma_{\rm piv}$ from
$q,M,\Delta_{\rm rnd},\Delta_{\rm aff}$ or other Pfaffian-format data.
